\documentclass[11pt]{article}
\usepackage{amsmath,amssymb,amsthm,mathtools,hyperref,enumitem,geometry}
\geometry{margin=1.1in}
\usepackage[T1]{fontenc}
\usepackage{microtype}
\usepackage{mathrsfs}
\usepackage{listings}
\lstset{basicstyle=\ttfamily\small,breaklines=true}

\title{Computing the Hodge Lattice of Fermat Hypersurfaces:\\
Rank and Discriminant via Shioda--Ran and Smith Normal Form}
\author{Prepared for a friendly duel with a persistent proposer}
\date{\today}

\newtheorem{theorem}{Theorem}
\newtheorem{lemma}{Lemma}
\newtheorem{prop}{Proposition}
\newtheorem{cor}{Corollary}
\theoremstyle{definition}
\newtheorem{definition}{Definition}
\newtheorem{remark}{Remark}

\newcommand{\Xmn}{X_m^n}
\newcommand{\PP}{\mathbf{P}}
\newcommand{\ZZ}{\mathbf{Z}}
\newcommand{\QQ}{\mathbf{Q}}
\newcommand{\CC}{\mathbf{C}}
\newcommand{\muone}{\mu_m^{\,n+2}/\mu_m}
\newcommand{\Hdg}{\mathrm{Hdg}}
\newcommand{\NS}{\mathrm{NS}}
\newcommand{\rank}{\mathrm{rank}}
\newcommand{\disc}{\mathrm{disc}}

\begin{document}
\maketitle

\begin{abstract}
We give a practical, computer-ready blueprint for computing the \emph{rank} and, when feasible, the \emph{discriminant} of the lattice of Hodge cycles in the middle cohomology of Fermat hypersurfaces
\(\Xmn : x_0^m+\cdots+x_{n+1}^m=0\subset \PP^{n+1}\) with \(n\) even.
Two complementary routes are described:
\begin{itemize}[nosep]
\item[(A)] a purely combinatorial count (after Shioda--Ran) giving \(\mathrm{rank}\) as \(1+\#\mathfrak{B}_m^{\,n}\);
\item[(B)] a Smith normal form (SNF) computation of a Gram matrix built either from (B1) linear algebraic cycles \( \PP^p \subset \Xmn\), or (B2) vanishing cycles projected to the Hodge subspace by character idempotents.
\end{itemize}
We also package ready-to-run \textsf{SageMath} code that implements (A) in full generality and (B1) for fourfolds. Worked examples include the Fermat surface of degree \(13\) and Fermat fourfolds of degrees \(5,7,9\).
\end{abstract}

\section{Set-up and the Shioda--Ran description}
Let \(n=2p\) be even and \(X=\Xmn\). Write \(G=\mu_m^{\,n+2}/\mu_m\) for the diagonal symmetry and \(\widehat G\cong \{(a_0,\dots,a_{n+1})\in (\ZZ/m\ZZ)^{n+2}\mid \sum a_i\equiv 0\}\) for its character group.
Following Shioda and Ran, the primitive middle cohomology decomposes into one-dimensional character eigenspaces \(V(\alpha)\) and the primitive Hodge classes are a direct sum of those \(V(\alpha)\) picked out by a simple congruence condition:

\begin{definition}[Shioda's set \(\mathfrak U_m^{\,n}\) and \(\mathfrak B_m^{\,n}\)]
Let \(\langle a\rangle\) denote the representative of \(a\in\ZZ/m\ZZ\) in \(\{1,\dots,m-1\}\). Put
\begin{align*}
\mathfrak U_m^{\,n}&=\Big\{\alpha=(a_0,\dots,a_{n+1})\in\widehat G\;:\; a_i\not\equiv 0\ \forall i\Big\},\\
\mathfrak B_m^{\,n}&=\Big\{\alpha\in\mathfrak U_m^{\,n}\;:\; \sum_{i=0}^{n+1}\big\langle t a_i\big\rangle =(p+1)m\ \text{ for all }t\in (\ZZ/m\ZZ)^\times\Big\}.
\end{align*}
\end{definition}

\begin{theorem}[Ran--Shioda]
If \(n=2p\) then the primitive Hodge classes are
\(\Hdg^p_{\mathrm{prim}}(X)=\bigoplus_{\alpha\in \mathfrak B_m^{\,n}} V(\alpha)\), each summand one-dimensional over \(\CC\). Consequently,
\begin{equation}\label{eq:rank-formula}
\rank\,\Hdg^p(X;\ZZ)=1+\#\,\mathfrak B_m^{\,n},
\end{equation}
the extra \(1\) being the non-primitive class \(h^p\) of the hyperplane.
\end{theorem}

\paragraph{Algorithm A (rank).} The condition defining \(\mathfrak B_m^{\,n}\) can be enforced with a small integer linear system on the \emph{multiplicity vector} \((c_1,\dots,c_{m-1})\), where \(c_r\) counts how many coordinates equal \(r\). Writing \(k=n+2\) and \(U=(\ZZ/m\ZZ)^\times\), one requires
\begin{equation}\label{eq:ILP}
\sum_{r=1}^{m-1} c_r = k,\qquad \sum_{r=1}^{m-1}\big\langle t r\big\rangle c_r = (p+1)m\quad(\forall\,t\in U),\qquad c_r\in\ZZ_{\ge 0}.
\end{equation}
For each solution \(c\), the number of ordered \((a_0,\dots,a_{k-1})\) it represents is the multinomial \(\binom{k}{c_1,\dots,c_{m-1}}\). Summing these multinomials over all solutions of \eqref{eq:ILP} yields \(\#\,\mathfrak B_m^{\,n}\), hence the rank by \eqref{eq:rank-formula}.

\medskip
\section{Intersection of linear cycles and a Gram matrix}
For \(n=2p\), the Fermat equation splits coordinates into \(p+1\) pairs; fixing a pairing and \(\zeta_j\in\mu_m\) per pair, the linear subspace
\begin{equation}\label{eq:linear-cycle}
L(\boldsymbol\zeta):\quad x_{2j}=\zeta_j\,x_{2j+1}\qquad (j=0,\dots,p)
\end{equation}
lies on \(\Xmn\). The intersection numbers inside \(X\) depend only on \(d=m\) and on the \emph{projective} dimension \(r=\dim(P_i\cap P_j)\) of their intersection \emph{in} \(\PP^{n+1}\):
\begin{equation}\label{eq:AMV-formula}
P_i\cdot P_j=\frac{1-(-d+1)^{\,r+1}}{d}\qquad(r=-1,0,1,\dots,p),
\end{equation}
with the convention \((\cdot)=0\) for $r=-1$.

\paragraph{Primitive part.} Pass to primitive classes \( \check P_i:=P_i-P_1\) and form the integer Gram matrix \(A_1=[\check P_i\cdot \check P_j]\). Its SNF gives the elementary divisors of the lattice generated by the linear cycles in primitive middle homology.

\paragraph{From primitive to full Hodge lattice.}
To pass from the primitive lattice to the full Hodge lattice one \emph{adjoins} the hyperplane class \(h^p\), which is orthogonal to the primitive part and satisfies \(\langle h^p,h^p\rangle = m\) in dimension \(n=4\).

\paragraph{Primitivity and discriminant.}
In favorable cases (e.g.\ degree \(5\) fourfolds) one can certify that the sublattice generated by linear cycles is \emph{primitive} in the Hodge lattice by comparing with the projector route below. When primitivity holds, the SNF yields the elementary divisors and hence the discriminant.

\end{document}

% --- APPENDIX: Projector/vanishing-cycle route (Algorithm B2) ---
\appendix
\section*{Appendix: Vanishing cycles and character projectors (Algorithm B2)}

For completeness we also package a projector-based computation that avoids algebraic cycles.
Consider the Brieskorn--Pham polynomial \(f=\sum x_i^m\). Its vanishing homology in degree \(n\) is
isomorphic to the primitive middle homology of the projective hypersurface \(X_m^n\).
A \(\mathbb{Z}\)-basis is naturally indexed by tuples \(u=(u_0,\dots,u_{n+1})\in\{1,\dots,m-1\}^{n+2}\).
The intersection form is the tensor product of \(A_{m-1}\)-forms, so on the tensor basis the
Gram matrix is simply the Kronecker product of the Cartan matrix; consequently the entry between
basis vectors \(u\) and \(v\) factors as a product of \(1\)-dimensional entries.

The diagonal symmetry \(G\) acts by characters \(\chi_u(g)=\zeta_m^{\langle u,t(g)\rangle}\). The
primitive \((p,p)\)-piece is the direct sum of eigenspaces with \(u\in\mathfrak{B}_m^n\). Hence
\emph{projecting} is equivalent to \emph{selecting} the basis vectors whose index \(u\) belongs to
\(\mathfrak{B}_m^n\). This gives an explicit integer Gram matrix for the primitive Hodge lattice.

The Sage file \texttt{fermat\_hodge\_B2\_projectors.sage} provides:
\begin{itemize}[nosep]
\item \texttt{list\_B\_tuples(m,n)} to enumerate \(\mathfrak{B}_m^n\) as tuples \(u\);
\item \texttt{primitive\_gram\_vanishing(m,n)} to assemble the Gram matrix on that sub-basis;
\item \texttt{full\_invariants\_from\_primitive} to adjoin the hyperplane class (\(\langle h^p,h^p\rangle=m\));
\item \texttt{compare\_with\_linear\_route\_fourfold(m)} to verify the equality with the linear-cycle route for \(n=4\).
\end{itemize}

\paragraph{Remark on normalization.}
For fourfolds (\(n=4\)) this construction matches the invariant factors reported by
Aljovin--Movasati--Villaflor for \(m=5\) when compared against the linear-cycle route. For \(m=7,9\),
running both routes and checking equality of Smith forms constitutes a computational proof that
the linear-cycle lattice is primitive in the full Hodge lattice, hence yields the discriminant.

